\section{Bayesian Evidence Fusion}

A key challenge in developing MERIT will be to address the uncertaity in conflicting signals that arise from handling multi-source candidate data.
MERIT employs a probabilistic framework based on Bayesian Evidence Fusion to address this challenge.
This is in contrast to traditional-rule based systems that rely on single sources of information.
By leveraging multiple sources, MERIT is able to provide a thorough "Glass-Box" assessment to separate the estimated ability from the certainty of that estimate.

\subsection*{The Beta Distribution Model}
We use a Beta distribution denoted as $Beta(\alpha, \beta)$ to model the "probability of competence" for a specific skill.
This distribution is well-suited for the task since its distribution can be updated incrementally as new evidence arrives.
This is key to providing an extensible, modular architecture that allows for the integration of new evidence sources beyond those that will be implemented in this project.

When we calculate the beta distribution for a metric, we can expect the following:
\begin{enumerate}
    \item \textbf{The Match Score:} The expected value (mean) of the distribution, providing the most likely estimate of a candidate's skill level.
    \item \textbf{The Confidence Level:} The spread (variance) of the distribution; a narrow distribution (low variance) signifies high confidence in the data, while a wide distribution (high variance) indicates a lack of evidence or conflicting signals.
\end{enumerate}

The parameters of the distribution are updated based on the type of evidence discovered:
\begin{itemize}
    \item \textbf{Alpha ($\alpha$):} Represents "success" signals or positive evidence (e.g., a verified GitHub contribution, a high-impact project on a CV).
    \item \textbf{Beta ($\beta$):} Represents "failure" signals, noise, or uncertainty (e.g., skill decay over time, lack of evidence for a self-reported claim).
\end{itemize}

The fused score (Match Score) is calculated as the mean of the distribution \cite{gelman2013bayesian}:
\begin{equation}
    E[X] = \frac{\alpha}{\alpha + \beta}
\end{equation}

The confidence level is calculated as the variance of the distribution:
\begin{equation}
    Var[X] = \frac{\alpha \beta}{(\alpha + \beta)^2 (\alpha + \beta + 1)}
\end{equation}

\subsection*{Modelling Uncertainty}
A key advantage of this framework is the ability to quantify \textit{uncertainty} through the variance of the distribution. 
As more evidence is gathered, the distribution narrows, increasing the confidence in the assessment. 
This allows MERIT to flag candidates with "High Confidence" vs. "Low Confidence" labels, providing the recruiter with a measure of how much to trust the automated ranking. 
By utilising the conjugate prior properties of the Beta distribution, MERIT can efficiently update its beliefs as new data sources are ingested.
